\subsection{Hilbertian spaces}

DEFINITION 5. --- A hilbertian space (or Hilbert space) is a prehilbertian space which is Hausdorff and complete. We say that a norm on a vector space E (over K) is hilbertian if it is prehilbertian, and if the normed space E is complete.

If E is a hilbertian space and M a closed vector subspace of E, the prehilbertian space structure induced on M is in fact a hilbertian space structure. In this case we say that M, with the induced structure is a hilbertian subspace of E.

Examples. --- 1) The prehilbertian spaces defined in examples 1, 3, 4 of V, p.~4, are hilbertian spaces. On the other hand, the prehilbertian space defined in example 2 is neither Hausdorff, nor complete. The complexification of a hilbertian space is a hilbertian space.

* 2) Let X be a Hausdorff topological space and let \(\mu\) be a positive measure on X. Let \(L^2(X, \mu)\) be the space consisting of equivalence classes, for \(\mu\) , of all square \(\mu\) -integrable functions on X with values in C. This is a complex hilbertian space, whose scalar product is given by

\[
\langle f | g \rangle = \int_ {\mathrm{X}} \overline {{{f (x)}}} g (x) d \mu (x). *
\]

* 3) Let \(n \geqslant 1\) be an integer and let \(U\) be an open set in \(\mathbf{R}^n\) . Let \(\mu\) be the measure on \(U\) induced by the Lebesgue measure on \(\mathbf{R}^n\) , and put \(\mathcal{H}^0 = L^2(U, \mu)\) . Let \(\mathcal{H}^1\) denote the space of all functions \(f \in \mathcal{H}^0\) with the following property; for \(1 \leqslant i \leqslant n\) , there exists a function \(g_i \in \mathcal{H}^0\) such that

\[
\int_ {\mathrm{U}} g _ {i} (x) h (x) d \mu (x) = - \int_ {\mathrm{U}} f (x) \mathrm{D} _ {i} h (x) d \mu (x)
\]

for every function \(h\) of class \(\mathbf{C}^1\) with compact support in \(\mathbf{U}\) . The function \(g_i\) is defined uniquely up to equivalence with respect to \(\mu\) , and is denoted by \(\mathbf{D}_i f\) or \(\partial f / \partial x_i\) (ith partial derivative). By induction on the integer \(s \geqslant 1\) , we define \(\mathcal{H}^s\) as the set of all functions \(f \in \mathcal{H}^1\) such that \(D_i f \in \mathcal{H}^{s-1}\) for \(1 \leqslant i \leqslant n\) . We define a scalar product on \(\mathcal{H}^s\) by the formula

\[
\langle f | g \rangle = \sum_ {k = 0} ^ {s} \sum_ {1 \leqslant i _ {1} \leqslant \dots \leqslant i _ {k} \leqslant n} \int \overline {{{{\mathrm{D} _ {i _ {1}} \dots \mathrm{D} _ {i _ {k}} f}}}}. \mathrm{D} _ {i _ {1}} \dots \mathrm{D} _ {i _ {k}} g d \mu .
\]

Then \(\mathcal{H}^s\) is a complex hilbertian space, called Sobolev space of index \(s\) .

* 4) Let X be a differential variety of class \(\mathbf{C}^r\) (with \(r \geqslant 1\) ) pure of finite dimension \(n\) .

In the vector fibre space \(\Lambda^{n}\mathrm{T}(\mathrm{X})\) , let L be the complement of the zero section. For every real number \(\lambda \neq 0\) , the mapping \(u \mapsto \lambda u\) from \(\Lambda^{n}\mathrm{T}(\mathrm{X})\) into itself leaves L stable.

Let \(\alpha\) be a complex number. A complex valued function \(\omega\) on \(\mathbf{L}\) such that \(\omega (\lambda u) = |\lambda |^{\alpha}\omega (u)\) for \(u\in \mathbf{L}\) and any non-zero real number \(\lambda\) is called a density of order \(\alpha\) on \(\mathbf{X}\) . We say that a density \(\omega\) of order 1 is locally integrable if there exists an open cover \((\mathrm{U}_i)_{i\in \mathbf{I}}\) of \(\mathbf{X}\) , and for every \(i\in \mathbf{I}\) a system of coordinates \(\xi_i = (\xi_i^1,\dots,\xi_i^n)\) on \(\mathrm{U}_i\) and a complex valued function \(f_{i}\) on \(\xi_i(\mathrm{U}_i)\) satisfying the following conditions: a) The function \(f_{i}\) is locally integrable on the open set \(\xi_i(\mathrm{U}_i)\) of \(\mathbf{R}^n\) with respect to the Lebesgue measure \(\mu\) ;

\begin{enumerate}
\def\labelenumi{\alph{enumi})}
\setcounter{enumi}{1}
\tightlist
\item
  Let \(x \in \mathbf{U}_i\) ; if \((\partial_{1,i,x}, \ldots, \partial_{n,i,x})\) is the basis of \(\mathrm{T}_x\mathrm{X}\) associated to the system of coordinates \((\xi_i^1, \ldots, \xi_i^n)\) in \(\mathbf{U}_i\) we have
\end{enumerate}

\[
\omega \left(\partial_ {1, i, x} \wedge \dots \wedge \partial_ {n, i, x}\right) = f _ {i} \left(\xi_ {i} ^ {1} (x), \dots , \xi_ {i} ^ {n} (x)\right).
\]

Then, there exists one and only one measure \(\tilde{\omega}\) on \(\mathbf{X}\) such that for every \(i\in \mathrm{I}\) , the image under \(\xi_{i}\) of the restriction of \(\tilde{\omega}\) to \(\mathrm{U}_i\) is equal to the measure \(f_{i}.\mu\) (cf.~VAR, R, 10.4.3).

Let \(\mathcal{V}\) (resp. \(\mathcal{N}\) ) be the vector space of measurable densities \(\omega\) of order \(1/2\) such that the measure associated with the density \(|\omega|^2\) of order 1 is bounded (resp. null). Let \(\omega_1\) and \(\omega_2\) be in \(\mathcal{V}\) ; then \(\omega = \overline{\omega}_1\omega_2\) is a density of order 1, and the measure \(\tilde{\omega}\) associated with \(\omega\) is bounded; the number \(\int_{\mathbf{X}}\tilde{\omega}\) depends only on the classes \(\dot{\omega}_1\) and \(\dot{\omega}_2\) of \(\omega_1\) and \(\omega_2\) modulo \(\mathcal{N}\) and is denoted by \(\langle\omega_1|\omega_2\rangle\) or \(\langle\dot{\omega}_1|\dot{\omega}_2\rangle\) . Then the mapping \((\dot{\omega}_1,\dot{\omega}_2)\mapsto\langle\dot{\omega}_1|\dot{\omega}_2\rangle\) assigns a complex hilbertian space structure to the vector space \(\Omega_{1/2}(\mathbf{X})=\mathcal{V}/\mathcal{N}_{*}\) .

*5) Let D be the open disc with centre 0 and radius 1 in C. The Hardy space \(\mathrm{H}^2 (\mathbf{D})\) consists of all holomorphic functions \(f\colon \mathbf{D}\to \mathbf{C}\) for which

\[
\sup _ {0 <   \mathbf {R} <   1} \int_ {0} ^ {1} | f (\mathbf {R}. \mathbf {e} (\theta)) | ^ {2} d \theta <   + \infty .
\]

If \(f_{1}\) and \(f_{2}\) belong to \(\mathrm{H}^2 (\mathbf{D})\) , the limit

\[
\langle f _ {1} | f _ {2} \rangle = \lim _ {\mathbf {R} \rightarrow 1} \int_ {0} ^ {1} \overline {{{f _ {1} (\mathbf {R} . \mathbf {e} (\theta))}}}. f _ {2} (\mathbf {R}. \mathbf {e} (\theta)) d \theta
\]

exists; the mapping \((f_1, f_2) \mapsto \langle f_1 | f_2 \rangle\) assigns a complex hilbertian space structure to the vector space \(\mathbf{H}^2(\mathbf{D})\) .

For a function \(f: \mathbf{D} \to \mathbf{C}\) to belong to \(\mathbf{H}^2(\mathbf{D})\) , it is necessary and sufficient that there exists a sequence \((a_n)_{n \in \mathbb{N}}\) of complex numbers such that \(\sum_{n=0}^{\infty} |a_n|^2 < +\infty\) and that

\[
f (z) = \sum_ {n = 0} ^ {\infty} a _ {n} z ^ {n}
\]

for all \(z \in \mathbf{D}\) . Then we have \(\| f\|^2 = \sum_{n=0}^{\infty} |a_n|^2\) which gives an isomorphism from \(\mathrm{H}^2(\mathrm{D})\) onto the hilbertian space \(\ell^2\) (V, p.~4). *

Every Hausdorff prehilbertian space is isomorphic to an everywhere dense subspace of a hilbertian space determined up to an isomorphism. Precisely :

PROPOSITION 4. --- Let E be a Hausdorff prehilbertian space, \(\hat{E}\) the normed space completion of E(GT, IX, § 3, No.~3). The scalar product \((x, y) \mapsto \langle x|y \rangle\) extends by continuity to a positive and separating hermitian form on \(\hat{\mathbf{E}}\) , and defines a hilbertian space structure on \(\hat{\mathbf{E}}\) .

The existence of the extension of \((x, y) \mapsto \langle x|y \rangle\) to \(\hat{E} \times \hat{E}\) follows from the continuity of this sesquilinear form on \(E \times E\) (GT, III, §6, No.~5, th. 1). Moreover, this extension, which will also be denoted by \((x, y) \mapsto \langle x|y \rangle\) is a hermitian form and satisfies the relation \(\langle x|x\rangle = \|x\|^2\) , by virtue of the principle of extension of identities ( \(\|x\|\) being the norm on \(\hat{E}\) obtained by extending the norm on E by continuity); this proves that the relation \(\langle x|x\rangle = 0\) implies x = 0 in \(\hat{E}\) , hence that the form \((x, y) \mapsto \langle x|y \rangle\) is positive and separating, and consequently defines a hilbertian space structure on \(\hat{E}\) .

Q.E.D.

This hilbertian space is said to be the completion of the Hausdorff prehilbertian space E.

* Example 6. --- Let U be an open subset of \(\mathbf{R}^n\) ( \(n \geqslant 1\) ). Let \(\mathcal{C}_0^1(\mathrm{U})\) be the vector space of all functions of class \(\mathbf{C}^1\) with compact support in U. We define a Hausdorff prehilbertian space structure on \(\mathcal{C}_0^1(\mathrm{U})\) whose scalar product is given by

\[
\langle f | g \rangle = \sum_ {i = 1} ^ {n} \int_ {\mathrm{U}} \overline {{{{\mathbf {D} _ {i} f (x)}}}}. \mathbf {D} _ {i} g (x) d x.
\]

This prehilbertian space is not complete. Its completion is called the Dirichlet space associated with U. *

COROLLARY. --- Let V be a vector space over K and f a positive hermitian form on V. a) There exists a Hilbert space E and a linear mapping \(u:V \rightarrow E\) such that \(f(x,y) = \langle u(x)|u(y) \rangle\) for x, y in V, and such that \(u(V)\) is dense in E.

\begin{enumerate}
\def\labelenumi{\alph{enumi})}
\setcounter{enumi}{1}
\tightlist
\item
  If two pairs \((\mathbf{E}_{i}, u_{i})\) satisfy the conditions analogous to a), then there exists a unique isomorphism \(\phi\) from the Hilbert space \(E_{1}\) onto the Hilbert space \(E_{2}\) such that \(u_{2} = \phi \circ u_{1}\) .
\end{enumerate}

Let N be the set of all \(x \in V\) such that \(f(x, x) = 0\) . We define a positive and separating hermitian form on the space V/N by \(\langle \dot{x} | \dot{y} \rangle = f(x, y)\) for \(x \in \dot{x}\) and \(y \in \dot{y}\) . Let E be the hilbertian space completion of V/N and u the mapping \(x \mapsto x + N\) from V into E. Then the conditions of a) are satisfied.

Under the hypotheses of b), N is equal to the kernel of \(u_{1}\) and to that of \(u_{2}\) . Hence there exists a bijective linear mapping \(\phi_{0}\) from \(u_{1}(V)\) onto \(u_{2}(V)\) such that \(u_{2}(x) = \phi_{0}(u_{1}(x))\) for all \(x \in V\) . We verify immediately that \(\phi_{0}\) is an isomorphism of prehilbertian spaces, hence an isometry. Since \(u_{i}(V)\) is dense in \(E_{i}\) for i = 1, 2, \(\phi_{0}\) extends uniquely to an isometry \(\phi\) from \(E_{1}\) onto \(E_{2}\) , and b) follows.

We say that the hilbertian space \(\mathbf{E}\) is the separated completion of \(\mathbf{V}\) (for the form \(f\) ).

Example 7. --- Let G be a group (with unit element 1) and \(\pi\) a homomorphism from G into the group of automorphisms of a complex hilbertian space E; we say that \(\pi\) is a unitary representation of G in E. Let \(a \in E\) ; we put

\[
\phi (x) = \langle a | \pi (x). a \rangle
\]

for all \(x \in \mathbf{G}\) . Then \(\phi : \mathbf{G} \to \mathbf{C}\) is positive definite, in other words satisfies the relation :

(PD) For every \(\lambda_{1},\ldots,\lambda_{n}\) in C and \(x_{1},\ldots,x_{n}\) in G, we have

\[
\sum_ {i, j = 1} ^ {n} \overline {{\lambda}} _ {i} \lambda_ {j} \phi (x _ {i} ^ {- 1} x _ {j}) \geqslant 0.\tag{11}
\]

In fact, the first member of (11) is precisely \(\| \sum_{i=1}^{n} \lambda_i \pi(x_i). a \|^2\) .

Conversely, let \(\phi\) be a positive definite function on \(G\) . Let \(\mathbf{C}^{(G)}\) be the vector space of all functions with finite support on \(G\) . We define a hermitian form \(\Phi\) on \(G\) by

\[
\Phi (u, v) = \sum_ {x, y \in G} \overline {{u (x)}} v (y) \phi (x ^ {- 1} y)\tag{12}
\]

and the relation (PD) expresses the fact that \(\Phi\) is positive. By the corollary of prop. 4, there exists a hilbertian space \(E\) and a linear mapping \(\rho : \mathbf{C}^{(G)} \to E\) , with a dense image, such that

\[
\Phi (u, v) = \langle \rho (u) | \rho (v) \rangle \quad \text { for } \quad u, v \quad \text { in } \quad \mathbf {C} ^ {(G)}.\tag{13}
\]

For every \(x \in \mathbf{G}\) , let \(\gamma_x\) be the left translation by \(x\) in \(\mathbf{C}^{(\mathrm{G})}\) defined by \(\gamma_x u(y) = u(x^{-1}y)\) for \(u \in \mathbf{C}^{(\mathrm{G})}\) and \(y \in \mathbf{G}\) . We have \(\Phi(\gamma_x u, \gamma_x v) = \Phi(u, v)\) . Now apply assertion \(b\) of the corollary of prop. 4 to \(\rho\) and \(\rho \circ \gamma_x\) : there exists a unique automorphism \(\pi(x)\) of the hilbertian space \(E\) such that \(\rho \circ \gamma_x = \pi(x) \circ \rho\) . We see immediately that \(\pi\) is a homomorphism from \(G\) into the group of automorphisms of \(E\) .

Let \(\delta\) be the element of \(\mathbf{C}^{(\mathrm{G})}\) defined by \(\delta(1) = 1\) , \(\delta(x) = 0\) for \(x \neq 1\) in \(\mathbf{G}\) . We have \(u = \sum_{x \in \mathbf{G}} u(x) \cdot \gamma_x \delta\) for all \(u \in \mathbf{C}^{(\mathrm{G})}\) , and so \(\rho(u) = \sum_{x \in \mathbf{G}} u(x) \pi(x) \cdot a\) by putting \(a = \rho(\delta)\) . Formulas (12) and (13) imply that \(\phi(x) = \langle a | \pi(x) \cdot a \rangle\) for all \(x \in \mathbf{G}\) . We remark that the set of vectors \(\pi(x) \cdot a\) for all \(x \in \mathbf{G}\) , is total in \(\mathbf{E}\) .

